\scanpage{66}
\mspage{- 10 -}
\margin{n° 384}
$\unsure{\mathcal M|\mathfrak M}_+(G)$ (ch. IV, § 1, n° 1, prop. 4) et l'ensemble des $\alpha\in T$ telles que $\alpha(U)\leq 1$ est donc fermé. Il en résulte que l'ensemble $\Delta_U$ des $\alpha\in T$ satisfaisant à $\alpha(\{e\})\geq 1$ et à $\alpha(U)\leq 1$ est \uline{fermé} : remarquons que ces conditions sont équivalentes à $\alpha(\{e\})=\alpha(U)=1$ et par suite, $\Delta_U$ se compose exactement de l'\uline{ensemble des mesures de Haar normalisées des sous-groupes discrets} $H$ \uline{tels que} $H\cap U=\{e\}$ ; nous noterons $D_U$ l'ensemble de ces sous-groupes discrets.

Soit maintenant $\varphi$ une fonction de $\mathcal K_+(G)$, de support contenu dans $U$ et telle que $\varphi(e)=1$ : on a alors $\alpha(\varphi)=1$ pour toute $\alpha\in\Delta_U$, ce qui montre que $\Delta_U$ est contenu dans le \uline{compact} $T\cap L_\varphi$ (prop. 2) et est donc lui-même \uline{compact}. L'application canonique de $\Delta_U$ sur $D_U$ est donc un homéomorphisme et on a :

PROPOSITION 7.- \uline{Pour tout voisinage ouvert} $U$ \uline{de} $e$, \uline{l'ensemble} $D_U$ \uline{des sous-groupes discrets} $H$ \uline{de} $G$ \uline{tels que} $H\cap U=\{e\}$ \uline{est compact}.

Soit maintenant $V$ un autre voisinage ouvert de $e$, relativement compact et tel que $\overline V\subset U$ et soit $\varphi\in\mathcal K_+(G)$, de support contenu dans $U$ et avec $\varphi(x)=1$ pour $x\in V$. On a $\alpha(\varphi)=1$ pour $\alpha\in\Delta_U$ et il existe donc un voisinage $W$ de $\alpha$ dans $\Delta$ tel que $\beta(\varphi)<2$ pour $\beta\in W$ : on a alors $W\subset\Delta_V$, ce qui montre que $\Delta_V$ est un voisinage de $\Delta_U$. Par suite, les \uline{intérieurs} des sous-espaces $\Delta_U$ de $\Delta$ forment un recouvrement ouvert de $\Delta$ par des ensembles relativement compacts. On en déduit (compte tenu de ce que $\Delta_{U_1}\cap\cdots\cap\Delta_{U_n}=\Delta_{U_1\cap\cdots\cap U_n}$) :

PROPOSITION 8.- \uline{Le sous-espace} $\Delta$ \uline{de} $T$ \uline{est localement compact ; pour qu'une partie} $A$ \uline{de} $\Delta$ \uline{soit relativement compacte, il faut et il suffit qu'il existe un voisinage ouvert} $U$ \uline{de} $e$ \uline{tel que} $A\subset\Delta_U$.

Supposons maintenant que $G$ satisfasse à la condition suivante :

(L)\qquad il existe un voisinage $Q$ de $e$ tel que le seul sous-groupe \uline{discret} de $G$ contenu dans $Q$ soit le sous-groupe réduit à $\{e\}$.

En termes imagés, la condition (L) signifie que $G$ n'a pas de sous-groupes discrets arbitrairement petits. * On peut montrer que tout groupe de Lie satis-
\margin{65}
